% Zbiór zadań — rozdział 4
% Rozdział jest określony przez nazwę tego pliku.

\begin{exo}[Całkowanie formy różniczkowej wzdłuż drogi]{integrale-forme-differentielle}
\keywords{forma różniczkowa, różniczka zupełna, kryterium Schwarza, całka krzywoliniowa}

Rozważmy dwie formy różniczkowe
\[\omega_1=y\,\mathrm{d}x+x\,\mathrm{d}y,\qquad\omega_2=y\,\mathrm{d}x,\]
oraz trzy drogi łączące początek układu $O(0,0)$ z punktem $M(1,1)$: drogę
$\gamma_1$, złożoną z poziomego odcinka $O\to(1,0)$, a następnie pionowego
odcinka $(1,0)\to M$; drogę $\gamma_2$, złożoną z pionowego odcinka
$O\to(0,1)$, a następnie poziomego odcinka $(0,1)\to M$; oraz drogę
$\gamma_3$, będącą odcinkiem prostej łączącym bezpośrednio $O$ z $M$.

\begin{enumerate}
\item Wykazać, że $\omega_1$ jest zupełna, znajdując funkcję $f$ taką, że $\omega_1=\mathrm{d}f$.
\item Obliczyć $\int_{\gamma_1}\omega_1$ i $\int_{\gamma_2}\omega_1$, parametryzując każdy odcinek. Otrzymać wynik również bez obliczeń.
\item Zastosować kryterium Schwarza do $\omega_2$. Jaki wniosek można wyciągnąć?
\item Obliczyć $\int_{\gamma_1}\omega_2$, $\int_{\gamma_2}\omega_2$ i $\int_{\gamma_3}\omega_2$. Skomentować wyniki.
\end{enumerate}

\begin{indication}
Na odcinku poziomym $\mathrm{d}y=0$, a na odcinku pionowym $\mathrm{d}x=0$:
każda całka sprowadza się do zwykłej całki pojedynczej.
\end{indication}

\begin{solution}
\textbf{Pytanie 1.} Odpowiednia jest funkcja $f(x,y)=xy$, ponieważ $\partial f/\partial x=y$ oraz $\partial f/\partial y=x$.

\textbf{Pytanie 2.} Dla krzywej sparametryzowanej przez $t\mapsto\bigl(x(t),y(t)\bigr)$ mamy
\[\int_\gamma\omega_1=\int\left[y(t)x'(t)+x(t)y'(t)\right]\,\mathrm{d}t.\]
Droga $\gamma_1$ jest sumą dwóch odcinków. Na pierwszym $x(t)=t$, $y(t)=0$
oraz $t\in[0,1]$, zatem $\mathrm{d}x=\mathrm{d}t$ i $\mathrm{d}y=0$. Na
drugim $x(t)=1$, $y(t)=t$, także dla $t\in[0,1]$, zatem $\mathrm{d}x=0$ i
$\mathrm{d}y=\mathrm{d}t$. W konsekwencji
\[\begin{aligned}I_1=\int_{\gamma_1}\omega_1&=\int_0^1(0\times1+t\times0)\,\mathrm{d}t+\int_0^1(t\times0+1\times1)\,\mathrm{d}t\\&=0+\int_0^1\mathrm{d}t=1.\end{aligned}\]
Dla $\gamma_2$ pierwszy odcinek parametryzujemy przez $x(t)=0$, $y(t)=t$,
a drugi przez $x(t)=t$, $y(t)=1$, w obu przypadkach dla $t\in[0,1]$. Otrzymujemy
\[\begin{aligned}I_2=\int_{\gamma_2}\omega_1&=\int_0^1(t\times0+0\times1)\,\mathrm{d}t+\int_0^1(1\times1+t\times0)\,\mathrm{d}t\\&=0+\int_0^1\mathrm{d}t=1.\end{aligned}\]
Obie wartości są równe. Można było przewidzieć to bez obliczeń: skoro
$\omega_1=\mathrm{d}f$ dla $f(x,y)=xy$, całka zależy wyłącznie od końców drogi,
\[\int_\gamma\omega_1=f(M)-f(O)=f(1,1)-f(0,0)=1.\]

\textbf{Pytanie 3.} Zapiszmy $\omega_2=A(x,y)\,\mathrm{d}x+B(x,y)\,\mathrm{d}y$.
Tutaj $A(x,y)=y$ i $B(x,y)=0$. Gdyby forma była zupełna, kryterium Schwarza dawałoby
\[\frac{\partial A}{\partial y}=\frac{\partial B}{\partial x}.\]
Tymczasem $\partial A/\partial y=1$, a $\partial B/\partial x=0$. Pochodne
mieszane są różne, więc $\omega_2$ nie jest zupełna. Jej całka może zatem
zależeć od drogi między $O$ i $M$, co potwierdza poniższe obliczenie.

\textbf{Pytanie 4.} Ponieważ $\omega_2=y\,\mathrm{d}x$, wkład do całki może
dać tylko zmiana $x$ zachodząca przy niezerowej wartości $y$. Korzystając z
poprzednich parametryzacji, na $\gamma_1$ otrzymujemy
\[\begin{aligned}J_1=\int_{\gamma_1}\omega_2&=\int_0^1 0\times1\,\mathrm{d}t+\int_0^1 t\times0\,\mathrm{d}t\\&=0.\end{aligned}\]
Na $\gamma_2$ pierwszy odcinek spełnia $x(t)=0$, więc $\mathrm{d}x=0$; na
drugim $x(t)=t$, $y(t)=1$ i $\mathrm{d}x=\mathrm{d}t$. Stąd
\[\begin{aligned}J_2=\int_{\gamma_2}\omega_2&=\int_0^1 t\times0\,\mathrm{d}t+\int_0^1 1\times1\,\mathrm{d}t\\&=1.\end{aligned}\]
Odcinek przekątny $\gamma_3$ ma natomiast parametryzację
\[x(t)=t,\qquad y(t)=t,\qquad t\in[0,1],\]
skąd $\mathrm{d}x=\mathrm{d}t$ i $\mathrm{d}y=\mathrm{d}t$. Zatem
\[J_3=\int_{\gamma_3}\omega_2=\int_0^1 y(t)x'(t)\,\mathrm{d}t=\int_0^1 t\,\mathrm{d}t=\frac12.\]
Trzy drogi mają te same końce, lecz dają trzy różne wartości: $J_1=0$,
$J_2=1$ i $J_3=1/2$. Jest to właśnie zależność od drogi charakterystyczna
dla formy różniczkowej, która nie jest zupełna.
\end{solution}
\end{exo}

\begin{exo}[Ta sama zmiana energii wewnętrznej, trzy sposoby przekazu]{meme-delta-u-trois-transferts}
\keywords{pierwsza zasada termodynamiki, energia wewnętrzna, ciepło, praca mechaniczna, praca elektryczna, kalorymetria, funkcja stanu}

Ciecz, kalorymetr i niewielki zanurzony opornik tworzą układ zamknięty,
makroskopowo pozostający w spoczynku. Naczynie jest sztywne, a całkowitą
pojemność cieplną układu przyjmujemy za stałą:
\[C=2{,}00\ \mathrm{kJ\,K^{-1}}.\]
Układ ma przejść ze stanu równowagi $A$ o temperaturze
$T_A=20{,}0\,^\circ\mathrm{C}$ do stanu równowagi $B$ o temperaturze
$T_B=25{,}0\,^\circ\mathrm{C}$. Przyjmujemy, że
\[U(B)-U(A)=C(T_B-T_A).\]
Zmiany makroskopowej energii kinetycznej i potencjalnej są równe zeru.
Pomijamy wszelkie straty do otoczenia.

Niezależnie realizujemy trzy poniższe protokoły, rozpoczynając od tego samego
stanu początkowego $A$:
\begin{itemize}
\item[Protokół 1] Układ zostaje doprowadzony do kontaktu cieplnego z gorącym źródłem, bez dostarczania mu pracy.
\item[Protokół 2] Ściany są izolowane cieplnie. Łopatka miesza ciecz dzięki opadaniu masy $m$ z wysokości $h=5{,}00$ m. Ostatecznie łopatka i ciecz spoczywają, a cała energia mechaniczna utracona przez masę zostaje przekazana układowi. Przyjąć $g=9{,}81\ \mathrm{m\,s^{-2}}$.
\item[Protokół 3] Układ otrzymuje ciepło $Q_3=4{,}00$ kJ, a pozostała energia jest dostarczana elektrycznie do opornika. Otrzymywana moc elektryczna jest stała i wynosi $\mathcal P=300$ W.
\end{itemize}

\begin{enumerate}
\item Obliczyć wspólną dla trzech protokołów zmianę energii wewnętrznej $\Delta U=U(B)-U(A)$.
\item Dla każdego z dwóch pierwszych protokołów wyznaczyć ciepło $Q$ i pracę $W$ otrzymane przez układ. W drugim obliczyć potrzebną masę $m$.
\item Dla trzeciego protokołu obliczyć otrzymaną pracę elektryczną i czas zasilania opornika.
\item Trzy przemiany mają te same stany początkowy i końcowy. Wyjaśnić, dlaczego ich wartości $Q$ i $W$ mogą być mimo to różne. Czy w stanie końcowym $B$ układ zawiera „ciepło” lub „pracę”?
\end{enumerate}

\begin{indication}
Zastosować $\Delta U=Q+W$ z konwencją bankiera. Praca otrzymana przez łopatkę
jest równa zmniejszeniu $mgh$ energii potencjalnej masy. Energię elektryczną
przechodzącą przez granicę przewodami liczy się jako pracę elektryczną.
\end{indication}

\begin{solution}
\textbf{Pytanie 1.} Różnica temperatur wynosi $T_B-T_A=5{,}0$ K. Wobec tego
\[\Delta U=C(T_B-T_A)=2{,}00\times5{,}0=10{,}0\ \mathrm{kJ}.\]
Wartość ta zależy wyłącznie od dwóch stanów równowagi $A$ i $B$.

\textbf{Pytanie 2.} W pierwszym protokole układ nie otrzymuje pracy: $W_1=0$.
Pierwsza zasada termodynamiki daje więc
\[Q_1=\Delta U=10{,}0\ \mathrm{kJ}.\]
Ciepło jest dodatnie, ponieważ energia wpływa do układu.

W drugim protokole ściany są izolowane cieplnie, zatem $Q_2=0$. Cała zmiana
energii wewnętrznej pochodzi od mechanicznej pracy łopatki:
\[W_2=\Delta U=10{,}0\ \mathrm{kJ}.\]
Ponieważ $W_2=mgh$, otrzymujemy
\[m=\frac{W_2}{gh}=\frac{10{,}0\times10^3}{9{,}81\times5{,}00}\simeq204\ \mathrm{kg}.\]
Duży rząd wielkości przypomina, że nawet niewielki wzrost temperatury
odpowiada już znacznej energii mechanicznej.

\textbf{Pytanie 3.} Pierwsza zasada wymaga
\[W_3=\Delta U-Q_3=10{,}0-4{,}00=6{,}00\ \mathrm{kJ}.\]
Praca ta jest dodatnia: zasilanie elektryczne dostarcza energię do układu.
Z $W_3=\mathcal P\,\Delta t$ wynika
\[\Delta t=\frac{W_3}{\mathcal P}=\frac{6{,}00\times10^3}{300}=20{,}0\ \mathrm{s}.\]

\textbf{Pytanie 4.} Trzy drogi dają odpowiednio
\[(Q_1,W_1)=(10{,}0,0)\ \mathrm{kJ},\qquad(Q_2,W_2)=(0,10{,}0)\ \mathrm{kJ},\]
oraz
\[(Q_3,W_3)=(4{,}00,6{,}00)\ \mathrm{kJ}.\]
W każdym przypadku suma $Q+W$ wynosi $10{,}0$ kJ i powoduje tę samą zmianę
$\Delta U$. Energia wewnętrzna jest funkcją stanu, natomiast ciepło i praca
opisują przekazy zachodzące podczas przemiany. W stanie końcowym $B$ układ ma
energię wewnętrzną, lecz nie zawiera ani „ciepła”, ani „pracy”.
\end{solution}
\end{exo}

\begin{exo}[Sprężanie pod stałym ciśnieniem zewnętrznym]{compression-pression-exterieure-constante}
\seoready{true}
\keywords{praca sił ciśnienia, stałe ciśnienie zewnętrzne, sprężanie, rozprężanie, rozprężanie swobodne}

Gaz jest sprężany od objętości $V_A$ do objętości $V_B<V_A$ pod stałym
ciśnieniem zewnętrznym $P_0$.

\begin{enumerate}
\item Obliczyć pracę otrzymaną przez gaz.
\item Gaz wraca od $V_B$ do $V_A$ przeciw temu samemu ciśnieniu zewnętrznemu $P_0$. Obliczyć otrzymaną pracę i porównać ją z pracą sprężania.
\item Ponownie rozważyć rozprężanie od $V_B$ do $V_A$, gdy gaz rozpręża się w próżni. Zinterpretować trzy otrzymane znaki.
\end{enumerate}

\begin{indication}
Zastosować $W=-\int P_{\mathrm{ext}}\,\mathrm{d}V$. Należy całkować ciśnienie
\emph{zewnętrzne}, również wtedy, gdy przemiana nie jest quasi-statyczna.
\end{indication}

\begin{solution}
\textbf{Pytanie 1.} Ciśnienie zewnętrzne jest stałe, zatem
\[W_{A\to B}=-P_0(V_B-V_A)=P_0(V_A-V_B)>0.\]
Podczas sprężania gaz otrzymuje pracę.

\textbf{Pytanie 2.} Dla drogi powrotnej
\[W_{B\to A}=-P_0(V_A-V_B)<0.\]
Obie prace są sobie przeciwne, ponieważ przemiany zachodzą przeciw temu samemu
stałemu ciśnieniu zewnętrznemu.

\textbf{Pytanie 3.} W próżni $P_{\mathrm{ext}}=0$, więc
\[W_{B\to A}^{\mathrm{vide}}=0.\]
Rozprężanie nie oznacza zatem automatycznie ujemnej pracy: gaz musi rzeczywiście
odpychać siłę zewnętrzną.
\end{solution}
\end{exo}
